% 原创教学示意：同一更新函数在不同步复用，状态不共享数值。
\begin{tikzpicture}[x={\dimexpr\linewidth/169\relax},y=1mm,
  font=\figurefont\footnotesize,text=NNInk]
  \node[nn operator,minimum width=22mm] (compact) at (18,0) {$F_{\bm\theta}$};
  \node[nn input] (xin) at (18,-23) {$\bm x_t$};
  \node[nn output] (yout) at (18,24) {$\widehat{\bm y}_t$};
  \draw[nn flow] (xin)--(compact);
  \draw[nn flow,draw=NNOutput] (compact)--(yout);
  \draw[nn operator path] (compact.east)--(35,0)--(35,12)--(18,12)--(compact.north);
  \node[nn annotation,anchor=west] at (27,18) {状态反馈};
  \draw[nn link] (45,-31)--(45,33);
  \node[nn hidden] (h0) at (57,0) {$\bm h^{(0)}$};
  \foreach \i/\xx in {1/84,2/117,3/150}{
    \node[nn operator,minimum width=20mm] (f\i) at (\xx,0) {$F_{\bm\theta}$};
    \node[nn input] (x\i) at (\xx,-23) {$\bm x_{\i}$};
    \node[nn output] (y\i) at (\xx,24) {$\widehat{\bm y}_{\i}$};
    \draw[nn flow] (x\i)--(f\i);
    \draw[nn flow,draw=NNOutput] (f\i)--(y\i);
  }
  \draw[nn operator path] (h0)--(f1);
  \draw[nn operator path] (f1)--node[above] {$\bm h^{(1)}$} (f2);
  \draw[nn operator path] (f2)--node[above] {$\bm h^{(2)}$} (f3);
  \draw[nn operator path] (f3.east)--(168,0);
  \node[anchor=south east] at (168,1) {$\bm h^{(3)}$};
  \draw[nn gradient] (145,-10)--(125,-10);
  \draw[nn gradient] (111,-10)--(92,-10);
\end{tikzpicture}
